\section{Strong Formulations}

In this section, we pass from the variational setting to stronger formulations of Signorini's problem. This is in many respects the central part of the report, since here the contact conditions become progressively more explicit. We first introduce suitable normal traces for stresses in $H(\diver;\Omega)$ and use them to derive a hybrid formulation. Based on this, we then recover the contact conditions in a KKT-type form, where nonpenetration, sign condition, and complementarity are clearly separated. Finally, under additional regularity assumptions, we arrive at a classical strong formulation in which these contact conditions attain their most transparent form.

\subsection{Stress Regularity and Normal Traces on \texorpdfstring{$H(\diver;\Omega)$}{Hdiv}}\label{sec:Hdiv}

The aim of this subsection is to develop a concept for normal traces of stress tensors $\sigma(u)$, in order to then subject them to contact conditions. Previously, in \eqref{eq:nonpenetration}, for $v\in H^1(\Omega;\R^d)$, we defined its classical normal trace $v_\mathrm{n}$ as the trace $\gamma v\in H^{1/2}(\Gamma;\R^d)$ times the normal $\mathrm{n}\in L^\infty(\Gamma;\R^d)$. Now, since $u\in H^1(\Omega;\R^d)$, the stress tensor $\sigma(u)$ only belongs to $L^2(\Omega;\S^d)$. Therefore, in general, a trace operator $\gamma$ is not available for $\sigma(u)$ and further regularity assumptions have to be taken on $\sigma(u)$, to get a notion for its trace. The canonical idea is to require $\sigma(u)\in H^1(\Omega;\S^d)$. In that case, $\gamma(\sigma(u))\in H^{1/2}(\Gamma;\S^d)$ which makes it possible to define a classical normal trace $\gamma(\sigma(u))\mathrm{n}\in H^{1/2}(\Gamma;\R^d)$. As $\sigma$ is an operator 
\begin{equation*}
	\sigma: H^s(\Omega;\R^d)\to H^{s-1}(\Omega;\S^d),\qquad \sigma(u)= \cA:\varepsilon(u),
\end{equation*}
this assumption further requires $u\in H^2(\Omega;\R^d)$ and also higher regularity on $\cA$, which is much stronger than needed but will also be assumed later on in Section~\ref{sec:classical}. A more relaxed assumption for now is $\sigma(u)\in H(\diver;\Omega)$, where
\begin{equation*}
	H(\diver;\Omega):=\left\{\tau\in L^2(\Omega;\S^d)\;\middle|\;\diver\tau\in L^2(\Omega;\R^d)\right\}
\end{equation*}
and to derive a concept for normal traces on $H(\diver;\Omega)$. Additionally to this assumption being a lot weaker, than $u\in H^2(\Omega;\R^d)$, it provides us with $\diver\sigma(u)\in L^2(\Omega;\R^d)$, which is desired in strong formulations as strong bulk equilibrium equations $-\diver\sigma(u) = f$ for source terms $f\in L^2(\Omega;\R^d)$ are now well-defined. It remains to provide a concept for normal traces on $H(\diver;\Omega)$. For sufficiently smooth $\tau:\overline{\Omega}\to\S^d$ and all sufficiently smooth $v:\overline{\Omega}\to\R^d$, {\em Green's identity} reads
\begin{equation}\label{eq:green_smooth}
	\int_\Gamma (\gamma\tau)\,\mathrm{n}\cdot(\gamma v)\;\ds  = \int_\Omega \tau : \nabla v\;\dx +\int_\Omega (\diver\tau)\cdot v\;\dx
\end{equation}
and can be interpreted as a definition of the normal trace $(\gamma\tau)\,\mathrm{n}$ in a weak sense, i.e.\ by integration against test functions. The key point now is, that the right-hand side of \eqref{eq:green_smooth} is still well-defined if $\tau$ and $v$ are not sufficiently smooth for the left-hand side to be well-defined, such as $\tau\in H(\diver;\Omega)$ and $v\in H^1(\Omega;\R^d)$, by understanding $\diver \tau$ and $\nabla v$ in the weak sense. Hence, for fixed $\tau\in H(\diver;\Omega)$, we define
\begin{equation}\label{eq:Ftau}
	F_\tau(v):= \int_\Omega \tau : \nabla v\;\dx +\int_\Omega (\diver\tau)\cdot v\;\dx,\qquad v\in H^1(\Omega;\R^d).
\end{equation}
By Cauchy--Schwarz,
\begin{equation*}
	|F_\tau(v)|
	\le
	\left(\|\tau\|_{0,\Omega}+\|\diver\tau\|_{0,\Omega}\right)\,\|v\|_{1,\Omega}
	\qquad\forall\,v\in H^1(\Omega;\R^d),
\end{equation*}
i.e.\ $F_\tau$ is a continuous linear functional on $H^1(\Omega;\R^d)$. Crucially, $F_\tau$ only depends on the trace $\gamma v\in H^{1/2}(\Gamma;\R^d)$ instead of the whole $v\in H^1(\Omega;\R^d)$.
\begin{proposition}[Dependence of $F_\tau$ on the Trace]\label{pr:Ftau_depends_on_trace}
	Let $\tau\in H(\diver;\Omega)$. Then, the functional $F_\tau$ defined in \eqref{eq:Ftau} vanishes on $H^1_0(\Omega;\R^d)$, i.e.
	\begin{equation*}
		F_\tau(w)=0
		\qquad\forall\,w\in H^1_0(\Omega;\R^d).
	\end{equation*}
	Consequently, $F_\tau$ depends only on the trace $\gamma v$ in the sense that
	\begin{equation*}
		\gamma v_1=\gamma v_2
		\qquad\Longrightarrow\qquad
		F_\tau(v_1)=F_\tau(v_2)
		\qquad\forall\,v_1,v_2\in H^1(\Omega;\R^d).
	\end{equation*}
\end{proposition}
\begin{proof}
	Let $w\in H^1_0(\Omega;\R^d)$. By density, there exists a sequence $(w_m)_{m\in\N}\subset \cD(\Omega;\R^d)$ such that
	\begin{equation*}
		w_m\to w
		\qquad\text{in }H^1(\Omega;\R^d).
	\end{equation*}
	Since $\tau\in H(\diver;\Omega)$, we have for every $m\in\N$ that
	\begin{equation*}
		\int_\Omega (\diver\tau)\cdot w_m\;\dx
		=
		-\int_\Omega \tau:\nabla w_m\;\dx,
	\end{equation*}
	and therefore
	\begin{equation*}
		F_\tau(w_m)=0
		\qquad\forall\,m\in\N.
	\end{equation*}
	Passing to the limit $m\to\infty$ and using the continuity of $F_\tau$ on $H^1(\Omega;\R^d)$ yields
	\begin{equation*}
		F_\tau(w)=\lim_{m\to\infty}F_\tau(w_m)=0.
	\end{equation*}
	Hence $F_\tau$ vanishes on $H^1_0(\Omega;\R^d)$. Now let $v_1,v_2\in H^1(\Omega;\R^d)$ satisfy $\gamma v_1=\gamma v_2$. Then
	\begin{equation*}
		\gamma(v_1-v_2)=0.
	\end{equation*}
	By the characterization $\ker(\gamma)=H^1_0(\Omega;\R^d)$, we obtain $v_1-v_2\in H^1_0(\Omega;\R^d)$. Therefore,
	\begin{equation*}
		F_\tau(v_1)-F_\tau(v_2)
		=
		F_\tau(v_1-v_2)
		=
		0,
	\end{equation*}
	which proves the claim.
\end{proof}
By Proposition~\ref{pr:Ftau_depends_on_trace}, the functional $F_\tau$ depends only on the trace $\gamma v$. Hence the map
\begin{equation*}
	\gamma v \mapsto F_\tau(v)
\end{equation*}
defines a well-defined continuous linear functional on $H^{1/2}(\Gamma;\R^d)$. The dual space of $H^{1/2}(\Gamma;\R^d)$ is denoted by $H^{-1/2}(\Gamma;\R^d)$ and the duality pairing by $\langle\cdot,\cdot\rangle_\Gamma$. By this definition, there exists a unique element of $H^{-1/2}(\Gamma;\R^d)$ representing the functional $F_\tau$.

\begin{definition}[Normal Trace on $H(\diver;\Omega)$]\label{def:normal_trace}
	Let $\tau\in H(\diver;\Omega)$. The {\tt normal trace} of $\tau$ is the unique element
	\begin{equation*}
		\gamma_{\mathrm{n}}(\tau)\in H^{-1/2}(\Gamma;\R^d)
	\end{equation*}
	such that
	\begin{equation}\label{eq:greenHdiv}
		\langle \gamma_{\mathrm{n}}(\tau),\,\gamma v\rangle_{\Gamma}
		=
		\int_\Omega \tau:\nabla v\;\dx + \int_\Omega (\diver\tau)\cdot v\;\dx
		\qquad\forall\,v\in H^1(\Omega;\R^d).
	\end{equation}
\end{definition}

In the smooth case, \eqref{eq:greenHdiv} reduces to \eqref{eq:green_smooth}, hence $\gamma_{\mathrm{n}}(\tau)$ coincides with the classical traction $(\gamma\tau)\,\mathrm n$. Coming back to our setting, we want to apply the derived theory to the symmetric stress tensor $\sigma(u)\in H(\diver;\Omega)$.

\begin{proposition}[Symmetric tensors act only on the symmetric gradient]\label{pr:symtensor}
	Let $\tau\in H(\diver;\Omega)$ and $v\in H^1(\Omega;\R^d)$. Then
	\begin{equation*}
		\tau:\nabla v=\tau:\varepsilon(v)
		\qquad \text{a.e.\ in }\Omega.
	\end{equation*}
\end{proposition}

\begin{proof}
	Since $\tau\in H(\diver;\Omega)\subset L^2(\Omega;\mathbb{S}^d)$, the tensor $\tau$ is symmetric a.e.\ in $\Omega$. Hence
	\begin{equation*}
		\tau:\nabla v
		=
		\tau:\varepsilon(v)
		+
		\tau:\left(\nabla v-\varepsilon(v)\right).
	\end{equation*}
	The term
	\begin{equation*}
		\nabla v-\varepsilon(v)=\frac12(\nabla v-\nabla v^\top)
	\end{equation*}
	is skew-symmetric, whereas $\tau$ is symmetric. Therefore their Frobenius product vanishes a.e., i.e.
	\begin{equation*}
		\tau:\left(\nabla v-\varepsilon(v)\right)=0
		\qquad \text{a.e.\ in }\Omega.
	\end{equation*}
	This proves
	\begin{equation*}
		\tau:\nabla v=\tau:\varepsilon(v)
		\qquad \text{a.e.\ in }\Omega.
	\end{equation*}
\end{proof}

By Proposition~\ref{pr:symtensor}, the normal trace $\sigma(u)\mathrm{n} := \gamma_\mathrm{n}(\sigma(u))\in H^{-1/2}(\Gamma;\R^d)$ of the symmetric stress tensor $\sigma(u)\in H(\diver;\Omega)$ is defined via
\begin{equation*}
	\langle\sigma(u)\mathrm{n},\,\gamma v\rangle_\Gamma = \int_\Omega\sigma(u):\varepsilon(v)\;\dx +\int_\Omega\diver\sigma(u)\cdot v\;\dx\qquad\forall\,v\in H^1(\Omega;\R^d),
\end{equation*}
according to \eqref{eq:greenHdiv}.

\subsection{A Hybrid Formulation}

The weak formulation in Problem \ref{pr:signorini2} encodes the equilibrium equation and the contact conditions only implicitly through the variational inequality. As discussed in the previous Section, under the additional stress regularity $\sigma(u)\in H(\diver;\Omega)$, these conditions can be separated into a bulk equation in $\Omega$ and boundary relations involving the traction $\sigma(u)\mathrm n$. This leads to the following hybrid formulation, which may be viewed as an intermediate form between the weak variational inequality and the strong Signorini conditions in the later subsections.

\begin{problem}[Signorini's Problem -- Hybrid Formulation]\label{pr:signorini4}
For given $f\in L^2(\Omega;\R^d)$ and $F\in L^2(\Gamma_N;\R^d)$, find $u\in K$ such that
\begin{subequations}\label{eq:hybrid}
	\begin{empheq}[left=\empheqlbrace]{alignat=2}
		-\diver\sigma(u) &= f
		&\qquad& \text{a.e.\ in }\Omega, \label{eq:hybrid_bulk}\\
		\langle \sigma(u)\mathrm n,\gamma u\rangle_\Gamma
		&= \mathrlap{\int_{\Gamma_N} F\cdot \gamma u\;\ds,}
		&& \label{eq:hybrid_eq}\\
		\langle \sigma(u)\mathrm n,\gamma v\rangle_\Gamma
		&\ge \int_{\Gamma_N} F\cdot \gamma v\;\ds
		&\qquad& \forall\, v\in K.
		\label{eq:hybrid_ineq}
	\end{empheq}
\end{subequations}
\end{problem}

We close this section by showing, that the hybrid formulation above is equivalent to the previous weak formulations in Problem \ref{pr:signorini2} in order to prove that it is indeed a valid formulation of the Signorini problem. For practical reasons in the proof, we start with the following identity.
\begin{lemma}
	Let $f\in L^2(\Omega;\R^d)$. If the bulk equilibrium equation $-\diver\sigma(u) = f$ holds almost everywhere in $\Omega$, then
	\begin{equation}\label{eq:useful}
		a(u,v)-\ell(v) = \langle\sigma(u)\mathrm{n}, \, \gamma v\rangle_\Gamma - \int_{\Gamma_N} F\cdot(\gamma v)\;\ds\qquad\forall\,v\in H^1(\Omega;\R^d).
	\end{equation}
\end{lemma}
\begin{proof}
	Using the bulk equation and the definition of $a(u,v)$, we have
	\begin{equation*}
		\langle\sigma(u)\mathrm{n}, \, \gamma v\rangle_\Gamma = \int_\Omega \sigma(u):\varepsilon(v)\;\dx +\int_\Omega\diver\sigma(u)\cdot v\;\dx = a(u,v) - \int_\Omega f\cdot v\;\dx\qquad \forall\,v\in H^1(\Omega;\R^d).
	\end{equation*}
	By the definition of $\ell$, this is
	\begin{equation*}
		a(u,v) - \int_\Omega f\cdot v\;dx = a(u,v)-\ell(v)+\int_{\Gamma_N} F\cdot(\gamma v)\;\ds \qquad \forall\,v\in H^1(\Omega;\R^d),
	\end{equation*}
	which proves the claim.
\end{proof}

\begin{proposition}[Equivalence of hybrid and weak formulation]
	If $u\in K$ solves Problem \ref{pr:signorini4}, then it also solves Problem \ref{pr:signorini2}. Conversely, if $u\in K$ solves Problem \ref{pr:signorini2} and if, in addition, $\sigma(u)\in H(\diver;\Omega)$, then it also solves Problem \ref{pr:signorini4}.
\end{proposition}
\begin{proof}
	Let $u\in K$ solve the Hybrid Problem \ref{pr:signorini4}. Subtracting \eqref{eq:hybrid_eq} from \eqref{eq:hybrid_ineq}, yields
	\begin{equation*}
		\langle \sigma(u)\mathrm{n},\,\gamma(v-u)\rangle_\Gamma - \int_{\Gamma_N} F\cdot \gamma(v-u)\;\ds\ge 0 \qquad\forall\,v\in K.
	\end{equation*}
	By \eqref{eq:hybrid_bulk}, we can directly apply the identity \eqref{eq:useful} with $v-u$ instead of $v$ to obtain
	\begin{equation*}
		a(u,v-u)-\ell(v-u)\ge 0 \qquad \forall\,v\in K,
	\end{equation*}
	which is exactly the weak formulation \ref{pr:signorini2}. Let now $u\in K$ solve the weak problem \ref{pr:signorini2} with the additional assumption that $\sigma(u)\in H(\diver;\Omega)$. Let $\varphi\in\cD(\Omega;\R^d)$ we a vector-valued test function on $\Omega$. Since $\varphi$ has compact support in $\Omega$ and thus a vanishing trace, $\gamma\varphi = 0$, we can perturb $u\in K$ by $\pm t\varphi,t\in \R$ without leaving $K$, i.e.\ $u\pm t\varphi\in K$. Testing both, $v = u+t\varphi$ and $v = u-t\varphi$, in the variational inequality of the weak formulation yields
	\begin{equation*}
		a(u,t\varphi)\ge \ell(t\varphi)\qquad\text{and}\qquad a(u,-t\varphi)\ge\ell(-t\varphi).
	\end{equation*}
	By linearity of $a$ in its second argument and linearity of $\ell$, we can pull out $t$ and $-t$ to obtain
	\begin{equation*}
		a(u,\varphi)\ge\ell(\varphi)\qquad\text{and}\qquad a(u,\varphi)\le\ell(\varphi)
	\end{equation*}
	and thus
	\begin{equation}\label{eq:phieq}
		a(u,\varphi) = \ell(\varphi)\qquad\forall\,\varphi\in\cD(\Omega;\R^d).
	\end{equation}
	The right-hand side of \eqref{eq:phieq} is
	\begin{equation*}
		\ell(\varphi) = \int_\Omega f\cdot\varphi\;\dx + \int_{\Gamma_N} F\cdot(\gamma\varphi)\;\ds = \int_\Omega f\cdot\varphi\;\dx,
	\end{equation*}
	where the boundary term vanished due to the compact support of $\varphi$. Furthermore, since $f\in L^2(\Omega;\R^d)$, it may be interpreted as a regular distribution, i.e.
	\begin{equation*}
		\int_\Omega f\cdot\varphi\;\dx = \langle f,\,\varphi\rangle, \qquad \forall\,\varphi\in\cD(\Omega;\R^d)
	\end{equation*}
	where $\langle\cdot,\,\cdot\rangle$ denotes the duality pairing between vector-valued distributions and vector-valued test functions. The right-hand side of \eqref{eq:phieq} is by Proposition~\ref{pr:symtensor}
	\begin{equation*}
		a(u,\varphi) = \int_\Omega\sigma(u):\varepsilon(\varphi)\;\dx = \int_\Omega\sigma(u):\nabla\varphi\;\dx.
	\end{equation*}
	Since $\sigma(u)\in H(\diver;\Omega)\subset L^2(\Omega;\S^d)$ it may also be seen as a regular distribution and the weak divergence is exactly defined by
	\begin{equation*}
		\langle -\diver\sigma(u),\,\varphi\rangle = \int_\Omega\sigma(u):\nabla \varphi\;\dx\qquad\forall\,\varphi\in\cD(\Omega;\R^d).
	\end{equation*}
	Thus, by \eqref{eq:phieq}, we have
	\begin{equation}\label{eq:weakbulk}
		\langle -\diver\sigma(u),\,\varphi\rangle = \langle f,\,\varphi\rangle\qquad\forall\,\varphi\in\cD(\Omega;\R^d).
	\end{equation}
	Now since $\sigma(u)\in H(\diver;\Omega)$, we also have $\diver\sigma(u)\in L^2(\Omega;\R^d)$ which means that \eqref{eq:weakbulk} reads
	\begin{equation*}
		-\diver\sigma(u) = f\qquad\text{a.e.\ in }\Omega,
	\end{equation*}
	which is exactly the equilibrium bulk equation \eqref{eq:hybrid_bulk}. Next, we prove the hybrid inequality \eqref{eq:hybrid_eq}. Since $K$ is a cone, $0\in K$ and $2u\in K$. Testing both, $v = 0$ and $v = 2u$, in the variational inequality of the weak formulation yields
	\begin{equation*}
		a(u,-u)\ge \ell(-u)\qquad\text{and}\qquad a(u,u)\ge \ell(u)
	\end{equation*}
	and thus $a(u,u) = \ell(u)$. Using the identity \eqref{eq:useful} with $u$ instead of $v$, we have
	\begin{equation*}
		0 = a(u,u)-\ell(u) = \langle \sigma(u)\mathrm{n},\,\gamma u\rangle_\Gamma-\int_{\Gamma_N}F\cdot\gamma u\;\ds,
	\end{equation*}
	which directly yields the hybrid equality \eqref{eq:hybrid_eq}. Moreover, with the identity \eqref{eq:useful} and the variational inequality from the weak formulation, we have
	\begin{equation*}
		\langle\sigma(u)\mathrm{n},\,\gamma(v-u)\rangle_\Gamma -\int_{\Gamma_N} F\cdot\gamma(v-u)\;\ds  = a(u,v-u)-\ell(v-u)\ge 0\qquad\forall\,v\in K.
	\end{equation*}
	By linearity, this is
	\begin{equation*}
		\langle\sigma(u)\mathrm{n},\,\gamma v\rangle_\Gamma\ge \int_{\Gamma_N}F\cdot \gamma v\;\ds + \left[\langle\sigma(u)\mathrm{n},\,\gamma u\rangle_\Gamma-\int_{\Gamma_N}F\cdot\gamma u\;\ds\right]\qquad \forall\,v\in V,
	\end{equation*}
	where the terms in the brackets cancel by the hybrid equality \eqref{eq:hybrid_eq}. Therefore, also the hybrid inequality \eqref{eq:hybrid_ineq} holds and the proof is completed.
\end{proof}

\subsection{Normal/Tangential Splitting and Local Normal Traces on \texorpdfstring{$H(\diver;\Omega)$}{Hdiv}}

In the hybrid formulation, the contact interaction is still encoded in the {\em global} normal trace
\begin{equation*}
	\sigma(u)\mathrm{n} = \gamma_{\mathrm n}(\sigma(u)) \in H^{-1/2}(\Gamma;\R^d),
\end{equation*}
which acts on traces on the whole boundary. We now restrict attention to the contact boundary $\Gamma_C$ and introduce suitable contact trace spaces together with the corresponding local contact traction $\gamma_{\mathrm{n},C}$. Throughout this subsection and the next subsection, we assume $|\Gamma_N| = 0$. Hence
\begin{equation*}
	\Gamma = \overline{\Gamma_D}\cup\overline{\Gamma_C},\qquad \Gamma_D\cap\Gamma_C = \emptyset
\end{equation*}
and every function in $V$ has trace supported only on $\Gamma_C$. We define the so-called {\em Lions-Magenes} space
\begin{equation*}
	H^{1/2}_{00}(\Gamma_C;\R^d)
	:=
	\left\{
	g\in H^{1/2}(\Gamma_C;\R^d)
	\;\middle|\;
	\mathrm{Ext}_0(g)\in H^{1/2}(\Gamma;\R^d)
	\right\},
\end{equation*}
where $\mathrm{Ext}_0(g)$ denotes the zero extension of $g$ from $\Gamma_C$ to $\Gamma$. Since $|\Gamma_N|=0$, the contact trace space is given by
\begin{equation*}
	W_C
	:=
	\left\{
	\gamma v|_{\Gamma_C}
	\;\middle|\;
	v\in V
	\right\}
	=
	H^{1/2}_{00}(\Gamma_C;\R^d).
\end{equation*}
Its norm is defined as the quotient norm
\begin{equation*}
	\|w\|_{W_C}
	:=
	\inf\left\{
	\|v\|_{1,\Omega}
	\;\middle|\;
	v\in V,\ \gamma v|_{\Gamma_C}=w
	\right\},
	\qquad w\in W_C.
\end{equation*}
We now split contact traces into normal and tangential parts. Define
\begin{equation*}
	W_{C,\mathrm n}
	:=
	\left\{
	w\cdot \mathrm{n}
	\;\middle|\;
	w\in W_C
	\right\},\qquad W_{C,t}
	:=
	\left\{
	w\in W_C
	\;\middle|\;
	w\cdot \mathrm{n} = 0
	\right\}.
\end{equation*}
Their quotient norms are given by
\begin{equation*}
	\|\eta\|_{W_{C,\mathrm n}}
	:=
	\inf\left\{
	\|v\|_{1,\Omega}
	\;\middle|\;
	v\in V,\ (\gamma v|_{\Gamma_C})\cdot \mathrm{n}=\eta
	\right\},
	\qquad \eta\in W_{C,\mathrm n},
\end{equation*}
and
\begin{equation*}
	\|\zeta\|_{W_{C,t}}
	:=
	\inf\left\{
	\|v\|_{1,\Omega}
	\;\middle|\;
	v\in V,\ (\gamma v|_{\Gamma_C})_t=\zeta
	\right\},
	\qquad \zeta\in W_{C,t}.
\end{equation*}
Every $w\in W_C$ admits the decomposition $w = w_{\mathrm n} \mathrm{n} + w_t$, with $w_{\mathrm n}\in W_{C,\mathrm n}$ and $w_t\in W_{C,t}$. Hence
\begin{equation*}
	W_C = W_{C,\mathrm n} \mathrm{n} \oplus W_{C,t}.
\end{equation*}
We denote the dual spaces of $W_C$, $W_{C,\mathrm n}$, and $W_{C,t}$ by $W_C^\prime,\, W_{C,\mathrm{n}}^\prime$ and $W_{C,t}^\prime$, respectively, and use the generic notation $\langle \cdot,\,\cdot\rangle_{\Gamma_C}$ for the corresponding duality pairings.

\begin{definition}[Contact normal trace]\label{def:normtang}
	Let $\tau\in H(\diver;\Omega)$. The contact normal trace of $\tau$ on $\Gamma_C$ is the functional
	\begin{equation*}
		\gamma_{\mathrm n,C}(\tau)\in W_C'
	\end{equation*}
	defined by
	\begin{equation*}
		\langle \gamma_{\mathrm n,C}(\tau),\,g\rangle_{\Gamma_C}
		:=
		\langle \gamma_{\mathrm n}(\tau),\,\mathrm{Ext}_0(g)\rangle_\Gamma
		\qquad
		\forall\, g\in W_C.
	\end{equation*}
\end{definition}

Since $W_C = H^{1/2}_{00}(\Gamma_C;\R^d)$ and $\mathrm{Ext}_0(g)\in H^{1/2}(\Gamma;\R^d)$ for every $g\in W_C$, this is well-defined. Moreover, if $g=\gamma v|_{\Gamma_C}$ for some $v\in V$, then $\gamma v=\mathrm{Ext}_0(g)$ on $\Gamma$ and accordingly
\begin{equation*}
	\langle \gamma_{\mathrm n,C}(\tau),g\rangle_{\Gamma_C}
	=
	\int_\Omega \tau:\nabla v\;\dx
	+
	\int_\Omega (\diver\tau)\cdot v\;\dx.
\end{equation*}
In particular, the right-hand side depends only on the contact trace $g=\gamma v|_{\Gamma_C}$ and not on the chosen lift $v\in V$, similar as in Proposition~\ref{pr:Ftau_depends_on_trace}.
\begin{definition}[Normal and tangential components of contact tractions]
	Let $\lambda\in W_C'$. Its normal component $\lambda_{C,\mathrm{n}}\in W_{C,\mathrm n}'$ and tangential component $\lambda_{C,t}\in W_{C,t}'$ are defined by
	\begin{equation*}
		\langle \lambda_{C,\mathrm{n}},\,\eta\rangle_{\Gamma_C}
		:=
		\langle \lambda,\,\eta \,\mathrm{n}\rangle_{\Gamma_C}
		\qquad
		\forall\,\eta\in W_{C,\mathrm n},
	\end{equation*}
	and
	\begin{equation*}
		\langle \lambda_{C,t},\,\zeta\rangle_{\Gamma_C}
		:=
		\langle \lambda,\,\zeta\rangle_{\Gamma_C}
		\qquad
		\forall\,\zeta\in W_{C,t}.
	\end{equation*}
\end{definition}
From now on, we denote by
\begin{equation*}
	\gamma_C v := (\gamma v)|_{\Gamma_C} \in W_C
\end{equation*}
the trace of $v \in V$ restricted to the contact boundary. Its normal and tangential components are denoted by
\begin{equation*}
	v_{\mathrm n} := \gamma_C v \cdot \mathrm n \in W_{C,\mathrm n},
	\qquad
	v_t := \gamma_C v - v_{\mathrm n}\mathrm n \in W_{C,t}.
\end{equation*}
With a slight abuse of notation, we use the same symbols $v_{\mathrm n}$ and $v_t$ for the restrictions to $\Gamma_C$ of the global normal and tangential components of $\gamma v$. Likewise, for $w \in W_C$ we write
\begin{equation*}
	w = w_{\mathrm n}\mathrm n + w_t,
	\qquad
	w_{\mathrm n} \in W_{C,\mathrm n}, \quad w_t \in W_{C,t}.
\end{equation*}
Moreover, we identify $W_{C,\mathrm n}'$ and $W_{C,t}'$ with the corresponding subspaces of $W_C'$, so that for $\lambda \in W_C'$ and $w = w_{\mathrm n}\mathrm n + w_t \in W_C$,
\begin{equation*}
	\langle \lambda, w \rangle_{\Gamma_C}
	=
	\langle \lambda_{C,\mathrm n}, w_{\mathrm n} \rangle_{\Gamma_C}
	+
	\langle \lambda_{C,t}, w_t \rangle_{\Gamma_C}.
\end{equation*}
Applying this construction to the stress tensor, we define the contact traction functional by
\begin{equation*}
	\sigma_C(u)\mathrm{n}
	:=
	\gamma_{\mathrm n,C}(\sigma(u))\in W_C'.
\end{equation*}
Its normal and tangential components are denoted by $\sigma_{C,\mathrm n}(u)\in W_{C,\mathrm n}'$ and $\sigma_{C,t}(u)\in W_{C,t}'$, respectively, such that
\begin{equation}\label{eq:sigmadecomp}
	\sigma_C(u)\mathrm{n} = \sigma_{C,\mathrm n}(u)\mathrm{n}+\sigma_{C,t}(u)\in W_C'.
\end{equation}
To encode the Signorini conditions, we introduce the cone of weakly nonpositive normal traces
\begin{equation*}
	W_{C,\mathrm n}^-
	:=
	\left\{
	\eta\in W_{C,\mathrm n}
	\;\middle|\;
	\eta\le 0 \text{ a.e.\ on }\Gamma_C
	\right\},
\end{equation*}
and its dual cone of weakly nonpositive normal contact forces
\begin{equation*}
	\Lambda_C
	:=
	\left\{
	\lambda\in W_{C,\mathrm n}'
	\;\middle|\;
	\langle \lambda,\,\eta\rangle_{\Gamma_C}\ge 0
	\;
	\forall\,\eta\in W_{C,\mathrm n}^-
	\right\}.
\end{equation*}
If $\lambda$ is represented by an $L^2$-function, then the condition $\lambda\in\Lambda_C$ is precisely the weak formulation of the condition $\lambda\le 0$ almost everywhere on $\Gamma_C$, i.e.\ of weakly nonpositive normal contact forces.

\subsection{KKT-Type Strong Formulation}

We now reformulate the contact conditions in a form that separates the normal and tangential contact reactions on $\Gamma_C$. This leads to a {\em KKT-type formulation} of Signorini's problem, where the contact conditions are expressed through {\em primal feasibility}, {\em dual feasibility}, and {\em complementarity}. Here, ``KKT'' refers to the {\em Karush--Kuhn--Tucker} conditions from constrained optimization. In the present setting, the nonpenetration condition $u_{\mathrm n}\le 0$ on $\Gamma_C$ plays the role of the primal inequality constraint, the condition $\sigma_{C,\mathrm n}(u)\in \Lambda_C$ encodes the weak nonpositivity of the normal contact force, and the complementarity relation $\langle \sigma_{C,\mathrm n}(u),\,u_{\mathrm n}\rangle_{\Gamma_C}=0$ expresses that contact pressure can occur only where the gap vanishes.

\begin{problem}[Signorini's Problem -- Strong Formulation with KKT Contact Conditions]\label{pr:signorini5}
For given $f\in L^2(\Omega;\R^d)$, find $u\in V$ such that
\begin{subequations}\label{eq:strong}
	\begin{empheq}[left=\empheqlbrace]{alignat=2}
		-\diver \sigma(u) &= f
		&\qquad& \text{a.e.\ in }\Omega, \label{eq:strong_bulk}\\
		\sigma_{C,t}(u) &= 0
		&\qquad& \text{in }W_{C,t}^\prime, \label{eq:strong_tangential}\\
		u_\mathrm{n} &\le 0
		&\qquad& \text{a.e.\ on }\Gamma_C, \label{eq:strong_nonpen}\\
		\sigma_{C,\mathrm{n}}(u) &\in \Lambda_C,
		&& \label{eq:strong_lambda}\\
		\langle\sigma_{C,\mathrm{n}}(u),\, u_\mathrm{n}\rangle_{\Gamma_C} &= 0.
		&& \label{eq:strong_comp}
	\end{empheq}
\end{subequations}
\end{problem}

Under the standing assumption $|\Gamma_N|=0$, this formulation is equivalent to the hybrid formulation introduced above. This equivalence will be proved now and will serve as the basis for the mixed formulation in Section~\ref{sec:mixed}.


\begin{proposition}[Equivalence of KKT and Hybrid Formulation]
	Suppose that $|\Gamma_N|= 0$. Then every solution to Problem \ref{pr:signorini5} is also a solution to Problem \ref{pr:signorini4} and vice versa.
\end{proposition}
\begin{proof}
	Assume first that $u\in V$ solves the above KKT formulation \ref{pr:signorini5} of Signorini's Problem. The bulk equation $-\diver\sigma(u) = f$ almost everywhere in $\Omega$ is contained in both formulations and by \eqref{eq:strong_nonpen}, $u\in K$. It remains to show the hybrid equality \eqref{eq:hybrid_eq} and inequality \eqref{eq:hybrid_ineq}, where in each the right-hand sides vanish, due to $|\Gamma_N| = 0$. Let $v\in K$ be arbitrary. Since $|\Gamma_N| = 0$ and $(\gamma v)|_{\Gamma_D} = 0$, we have
	\begin{equation*}
		\langle \sigma(u)\mathrm n,\,\gamma v\rangle_\Gamma = \langle \sigma_C(u)\mathrm n,\,\gamma_C v\rangle_{\Gamma_C}.
	\end{equation*}
	Using the decomposition \eqref{eq:sigmadecomp} of $\sigma_C(u)\mathrm{n}$ into normal and tangential components together with Definition~\ref{def:normtang}, we get
	\begin{equation*}
		\langle \sigma_C(u)\mathrm n,\,\gamma v\rangle_{\Gamma_C} = \langle \sigma_{C,\mathrm n}(u),\,v_\mathrm n\rangle_{\Gamma_C} + \langle \sigma_{C,t}(u),\,v_t\rangle_{\Gamma_C}.
	\end{equation*}
	By \eqref{eq:strong_tangential}, the second term vanishes. Therefore, $\langle \sigma(u)\mathrm n,\,\gamma v\rangle_\Gamma=\langle \sigma_{C,\mathrm{n}}(u),\,v_\mathrm n\rangle_{\Gamma_C}$. Since $v\in K$, we have $v_\mathrm{n}\le 0$ almost everywhere on $\Gamma_C$, and by \eqref{eq:strong_lambda} we know that $\sigma_{C,\mathrm{n}}(u)\in\Lambda_C$. Hence
	\begin{equation*}
		\langle \sigma_C(u)\mathrm n,\,\gamma_C v\rangle_{\Gamma_C}=\langle \sigma_{C,\mathrm{n}}(u),\,v_\mathrm n\rangle_{\Gamma_C}\ge 0,
	\end{equation*}
	which proves the inequality \eqref{eq:hybrid_ineq}. Taking $v=u$ and using \eqref{eq:strong_comp}, we similarly obtain
	\begin{equation*}
		\langle \sigma(u)\mathrm n,\,\gamma u\rangle_\Gamma
		=
		\langle \sigma_{C,\mathrm{n}}(u),\,u_\mathrm n\rangle_{\Gamma_C}
		=0,
	\end{equation*}
	which is exactly the hybrid equality \eqref{eq:hybrid_eq}. Thus $u$ solves Problem \ref{pr:signorini4}. Conversely, assume that $u\in K$ solves Problem \ref{pr:signorini4}. Then $u\in V$, and \eqref{eq:strong_nonpen} is already contained in the definition of $K$. Moreover, \eqref{eq:hybrid_bulk} is exactly \eqref{eq:strong_bulk}. We next prove \eqref{eq:strong_tangential}. Let $w_t\in W_{C,t}$. Then there exists $w\in V$ such that
	\begin{equation*}
		\gamma w|_{\Gamma_C}=w_t
		\qquad\text{and}\qquad
		w_\mathrm n=0\text{ on }\Gamma_C.
	\end{equation*}
	Since $u\in K$ and $w_\mathrm n=0$ on $\Gamma_C$, we have $u\pm w\in K$. Hence, the hybrid inequality \eqref{eq:hybrid_ineq} applied to $v=u+w$ and $v=u-w$, together with the hybrid equality \eqref{eq:hybrid_eq}, yields
	\begin{equation*}
		\langle \sigma(u)\mathrm n,\,\gamma w\rangle_\Gamma\ge 0
		\qquad\text{and}\qquad
		\langle \sigma(u)\mathrm n,\,\gamma w\rangle_\Gamma\le 0.
	\end{equation*}
	Therefore, $\langle \sigma(u)\mathrm n,\,\gamma w\rangle_\Gamma=0$. Again using $|\Gamma_N| = 0$ and $\gamma w = 0$ on $\Gamma_D$, we obtain
	\begin{equation*}
		0 = \langle\sigma(u)\mathrm{n},\,\gamma w\rangle_\Gamma = \langle\sigma_C(u)\mathrm{n},\,\gamma_C w\rangle_{\Gamma_C} = \langle \sigma_{C,\mathrm{n}}(u),\,w_\mathrm n\rangle_{\Gamma_C}+\langle \sigma_{C,t}(u),\,w_t\rangle_{\Gamma_C}=\langle \sigma_{C,t}(u),\,w_t\rangle_{\Gamma_C},
	\end{equation*}
	since $w_\mathrm{n} = 0$ on $\Gamma_C$. Since $w_t\in W_{C,t}$ was arbitrary, this proves \eqref{eq:strong_tangential}. We now prove \eqref{eq:strong_lambda}. Let $\eta\in W_{C,\mathrm{n}}^{-}$, i.e.\ $\eta\le 0$ a.e.\ on $\Gamma_C$. Then there exists $v\in V$ such that $v_\mathrm n=\eta$. Because $\eta\le 0$, we have $v\in K$. Hence the hybrid inequality \eqref{eq:hybrid_ineq} gives $\langle \sigma(u)\mathrm n,\,\gamma v\rangle_\Gamma\ge 0$. Using $|\Gamma_N|=0$, $\gamma v=0$ on $\Gamma_D$, and \eqref{eq:strong_tangential}, we infer
	\begin{equation*}
		0\le \langle \sigma_C(u)\mathrm n,\gamma_C v\rangle_{\Gamma_C}
		=\langle \sigma_{C,\mathrm{n}}(u),v_\mathrm n\rangle_{\Gamma_C}
		+\langle \sigma_{C,t}(u),v_t\rangle_{\Gamma_C}
		=\langle \sigma_{C,\mathrm{n}}(u),\eta\rangle_{\Gamma_C},
	\end{equation*}
	which proves $\sigma_{C,\mathrm{n}}(u)\in\Lambda_C$, i.e.\ \eqref{eq:strong_lambda}. Finally, the hybrid equality \eqref{eq:hybrid_eq}, together with $|\Gamma_N|=0$, $\gamma u=0$ on $\Gamma_D$, and \eqref{eq:strong_tangential}, yields
	\begin{equation*}
		0=\langle \sigma(u)\mathrm n,\gamma u\rangle_\Gamma
		=\langle \sigma_C(u)\mathrm n,\gamma_C u\rangle_{\Gamma_C}
		=\langle \sigma_{C,\mathrm{n}}(u),u_\mathrm n\rangle_{\Gamma_C}
		+\langle \sigma_{C,t}(u),u_t\rangle_{\Gamma_C}
		=\langle \sigma_{C,\mathrm{n}}(u),u_\mathrm n\rangle_{\Gamma_C},
	\end{equation*}
	which is exactly \eqref{eq:strong_comp}. Hence $u$ solves Problem \ref{pr:signorini5}.
\end{proof}

\subsection{Classical Strong Formulation under Additional Regularity}\label{sec:classical}

So far, the displacement field was sought only in $u\in V\subset H^1(\Omega;\R^d)$. Accordingly, the stress tensor satisfies only $\sigma(u)\in L^2(\Omega;\mathbb S^d)$ or, under the additional assumption discussed in Section~\ref{sec:Hdiv}, $\sigma(u)\in H(\diver;\Omega)$, so that no classical boundary trace of $\sigma(u)$ is available in general but only the weak concept of normal traces $\gamma_{\mathrm{n}}(\sigma(u))\in H^{-1/2}(\Gamma;\R^d)$ on $H(\diver;\Omega)$. In order to formulate the boundary traction and the contact conditions pointwise on $\Gamma_N$ and $\Gamma_C$, we now impose additional regularity on the displacement. We assume that
\begin{equation*}
	u\in H^s(\Omega;\R^d)
	\qquad\text{with}\qquad
	s>\frac32.
\end{equation*}
According to Definition~\ref{def:smallstraintensor}, the strain $\varepsilon(u)$ imposed by $u\in H^s(\Omega;\R^d)$ is now a function in $H^{s-1}(\Omega;\S^d)$. We further assume that the elasticity tensor field $\mathcal A:\Omega\to\mathcal L(\mathbb S^d;\mathbb S^d)$ is sufficiently regular so that
\begin{equation*}
	\sigma(u)=\mathcal A:\varepsilon(u)\in H^{s-1}(\Omega;\mathbb S^d).
\end{equation*}
For instance, this is the case if $\mathcal A$ is constant. More generally, for $\frac32<s\le 2$, it is sufficient to assume
\begin{equation*}
	\mathcal A_{ijkl}\in W^{1,\infty}(\Omega)
	\qquad\forall\,1\le i,j,k,l\le d.
\end{equation*}
In order to apply the trace theorem to $\sigma(u)$, we additionally assume
\begin{equation*}
	\Gamma\in \mathcal C^{m,1}
	\qquad\text{with}\qquad
	m\ge \lceil s-1\rceil.
\end{equation*}
Then $\gamma(\sigma(u))\in H^{s-3/2}(\Gamma;\mathbb S^d)$, and, since now the unit outward normal satisfies $\mathrm{n}\in \mathcal C^{m-1,1}(\Gamma;\R^d)$, the classical boundary traction
\begin{equation*}
	t(u):=(\gamma\sigma(u))\,\mathrm{n}
\end{equation*}
is well-defined and belongs to $H^{s-3/2}(\Gamma;\R^d)$. Its normal and tangential components are given by
\begin{equation*}
	\sigma_{\mathrm n}(u):=t(u)\cdot \mathrm{n}\in H^{s-3/2}(\Gamma),
	\qquad
	\sigma_t(u):=t(u)-\sigma_{\mathrm n}(u)\,\mathrm{n}\in H^{s-3/2}(\Gamma;\R^d).
\end{equation*}
Likewise, the displacement trace satisfies $\gamma u\in H^{s-1/2}(\Gamma;\R^d)$, with normal and tangential components
\begin{equation*}
	u_{\mathrm n}:=(\gamma u)\cdot \mathrm{n}\in H^{s-1/2}(\Gamma),
	\qquad
	u_t:=\gamma u-u_{\mathrm n}\,\mathrm{n}\in H^{s-1/2}(\Gamma;\R^d).
\end{equation*}
If, in addition, $\sigma(u)\in H(\diver;\Omega)$, then the normal trace on $H(\diver;\Omega)$ and the classical normal trace on $H^s(\Omega;\S^d)$ coincide in the sense that
\begin{equation}\label{eq:coincide}
	\langle \sigma(u)\mathrm{n},\gamma v\rangle_\Gamma
	=
	\int_\Gamma t(u)\cdot \gamma v\,\mathrm ds
	\qquad\forall\,v\in V.
\end{equation}
This allows us to pass from the hybrid formulation to a classical strong formulation with pointwise contact conditions.


\begin{problem}[Signorini's Problem -- Classical Strong Formulation]\label{pr:signorini6}
For given $f\in L^2(\Omega;\R^d)$ and $F\in L^2(\Gamma_N;\R^d)$ find $u\in H^s(\Omega;\R^d)$ with $s>3/2$ such that the {\em small strain elasticity equations}
\begin{subequations}\label{eq:strong_classical}
	\begin{empheq}[left=\empheqlbrace]{alignat=2}
		-\diver\sigma(u) &= f
		&\qquad& \text{in }\Omega, \label{eq:strong_classical_bulk}\\
		u &= 0
		&\qquad& \text{on }\Gamma_D, \label{eq:strong_classical_dirichlet}\\
		t(u) &= F
		&\qquad& \text{on }\Gamma_N, \label{eq:strong_classical_neumann}
	\end{empheq}
\end{subequations}
and the {\em contact conditions}
\begin{subequations}\label{eq:strong_contact}
	\begin{empheq}[left=\empheqlbrace]{alignat=2}
		u_\mathrm{n} &\le 0
		&\qquad& \text{on }\Gamma_C, \label{eq:strong_contact_nonpen}\\
		\sigma_\mathrm{n}(u) &\le 0
		&\qquad& \text{on }\Gamma_C, \label{eq:strong_contact_sigma_n}\\
		\sigma_\mathrm{n}(u)\,u_\mathrm{n} &= 0
		&\qquad& \text{on }\Gamma_C, \label{eq:strong_contact_comp}\\
		\sigma_t(u) &= 0
		&\qquad& \text{on }\Gamma_C \label{eq:strong_contact_tangential}
	\end{empheq}
\end{subequations}
hold.
\end{problem}

The classical strong formulation admits a direct mechanical interpretation. The equilibrium equation \eqref{eq:strong_classical_bulk} expresses, just as previously, the balance of internal stresses and body forces in the bulk. The homogeneous Dirichlet condition \eqref{eq:strong_classical_dirichlet} on $\Gamma_D$ means, just as before, that the body is clamped along $\Gamma_D$, whereas the Neumann condition \eqref{eq:strong_classical_neumann} prescribes the boundary traction acting on the free part of the boundary. On the potential contact boundary $\Gamma_C$, the condition $u_{\mathrm n}\le 0$ is exactly the {\em nonpenetration} constraint, previously encoded in $K$, while $\sigma_{\mathrm n}(u)\le 0$ states that the foundation can exert only {\em compressive normal forces}. The {\em complementarity relation} $\sigma_{\mathrm n}(u)\,u_{\mathrm n}=0$ connects these two conditions. Either the body is detached from the foundation, in which case no normal contact pressure acts, or the body is in contact, in which case the normal gap vanishes. Finally, the {\em tangential condition} $\sigma_t(u)=0$ describes frictionless contact, meaning that the foundation does not exert any tangential force on the body.

\begin{proposition}[Equivalence of Classical Strong and Hybrid Formulation]
	Let $u\in H^s(\Omega;\R^d)$ with $s>3/2$ solve Problem \ref{pr:signorini6}. Then, it also solves Problem \ref{pr:signorini4}. Conversely, if $u\in K$ solves Problem \ref{pr:signorini4} and if, in addition, $u\in H^s(\Omega;\R^d)$ with $s>3/2$, then it also solves Problem \ref{pr:signorini6}.
\end{proposition}
\begin{proof}
	Again, the bulk equation $-\diver\sigma(u) = f$ almost everywhere in $\Omega$ appears in both formulations. Let $u\in H^s(\Omega;\R^d),\, s>3/2$ solve Problem \ref{pr:signorini6}. With \eqref{eq:strong_classical_dirichlet} and \eqref{eq:strong_contact_nonpen} we have $u\in K$. It remains to prove the hybrid inequality \eqref{eq:hybrid_ineq} and equality \eqref{eq:hybrid_eq}, which here read
	\begin{equation}\label{eq:hybineq}
		\int_\Gamma t(u)\cdot(\gamma v)\;\ds \ge \int_{\Gamma_N}F\cdot(\gamma v)\;\ds\qquad \forall\,v\in K
	\end{equation}
	and
	\begin{equation}\label{eq:hybeq}
		\int_\Gamma t(u)\cdot(\gamma u)\;\ds = \int_{\Gamma_N} F\cdot(\gamma u)\;\ds,
	\end{equation}
	due to the sufficient regularity of $u$ such that the normal trace $\sigma(u)\mathrm{n}$ and the classical traction $t(u)$ coincide in the sense of \eqref{eq:coincide}. We decompose
	\begin{equation}\label{eq:decompint}
		\int_\Gamma t(u)\cdot(\gamma v) \;\ds = \int_{\Gamma_D}t(u)\cdot(\gamma v) \;\ds+\int_{\Gamma_N}t(u)\cdot(\gamma v) \;\ds+\int_{\Gamma_C}t(u)\cdot(\gamma v) \;\ds.
	\end{equation}
	By $(\gamma v)|_{\Gamma_D} = 0$ and the Neumann condition \eqref{eq:strong_classical_neumann}, we are left over with
	\begin{equation*}
		\int_\Gamma t(u)\cdot(\gamma v)\;\ds = \int_{\Gamma_N}F\cdot(\gamma v)\;\ds + \int_{\Gamma_C} t(u)\cdot(\gamma v)\;\ds.
	\end{equation*}
	With the splittings $t(u) = \sigma_\mathrm{n}(u)\mathrm{n}+\sigma_t(u)$ and $\gamma v = v_\mathrm{n}\mathrm{n}+v_t$ in their normal and tangential components, the integrand of the right term is, on $\Gamma_C$, given by
	\begin{equation*}
		t(u) \cdot(\gamma v) = \sigma_\mathrm{n}(u)\,v_\mathrm{n} + \sigma_t(u)\cdot v_t = \sigma_\mathrm{n}(u)\,v_\mathrm{n},
	\end{equation*}
	where the tangential term vanishes on $\Gamma_C$ due to \eqref{eq:strong_contact_tangential}. Furthermore, since $v\in K$, we have $v_\mathrm{n}\le 0$ on $\Gamma_C$. Together with \eqref{eq:strong_contact_sigma_n}, it follows that $\sigma_\mathrm{n}(u)v_\mathrm{n}\ge 0$ on $\Gamma_C$, i.e.
	\begin{equation*}
		\int_\Gamma t(u)\cdot(\gamma v)\;\ds = \int_{\Gamma_N}F\cdot(\gamma v)\;\ds + \int_{\Gamma_C}\sigma_\mathrm{n}(u)\,v_\mathrm{n}\;\ds \ge \int_{\Gamma_N}F\cdot(\gamma v)\;\ds\qquad\forall\, v\in K,
	\end{equation*}
	which is exactly the hybrid inequality \eqref{eq:hybineq}. By the same arguments with $v = u$, we arrive at
	\begin{equation*}
		\int_\Gamma t(u)\cdot(\gamma u)\;\ds = \int_{\Gamma_N}F\cdot(\gamma u)\;\ds +\int_{\Gamma_C}\sigma_\mathrm{n}(u)\,u_{\mathrm{n}} = \int_{\Gamma_N}F\cdot(\gamma u)\;\ds,
	\end{equation*}
	where the $\Gamma_C$-term vanishes due to the complementarity condition \eqref{eq:strong_contact_comp}. This proves the hybrid equality \eqref{eq:hybeq}. Therefore, $u$ is also a solution to Problem \ref{pr:signorini4}. Now, let $u\in K\cap H^s(\Omega;\R^d)$ with $s>3/2$ solve Problem \ref{pr:signorini4}. The requirement $u\in K$ immediately implies \eqref{eq:strong_classical_dirichlet} and \eqref{eq:strong_contact_nonpen}. We now prove the Neumann condition $t(u) = F$ on $\Gamma_N$, i.e.\ \eqref{eq:strong_classical_neumann}. For some $g\in H_{00}^{1/2}(\Gamma_N;\R^d)$ choose $v\in H^1(\Omega;\R^d)$ such that
	\begin{equation*}
		(\gamma v)|_{\Gamma_N} = g,\qquad (\gamma v)|_{\Gamma_C} = 0,\qquad (\gamma v)|_{\Gamma_D} = 0.
	\end{equation*}
	Since $v_\mathrm{n}= 0$ on $\Gamma_C$ and $\gamma v = 0$ on $\Gamma_D$, we have $v\in K$ and also $-v\in K$. Thus, we can test $\pm v$ in the hybrid inequality \eqref{eq:hybineq} to obtain
	\begin{equation*}
		\int_{\Gamma_N}t(u)\cdot g\;\ds\ge \int_{\Gamma_N} F\cdot g\;\ds,\qquad\text{and}\qquad -\int_{\Gamma_N}t(u)\cdot g\;\ds\ge -\int_{\Gamma_N} F\cdot g\;\ds
	\end{equation*}
	for all $g\in H_{00}^{1/2}(\Gamma_N;\R^d)$. Therefore,
	\begin{equation*}
		\int_{\Gamma_N} (t(u)- F)\cdot g\;\ds = 0\qquad\forall\,g\in H_{00}^{1/2}(\Gamma_N;\R^d),
	\end{equation*}
	i.e.\ $t(u) = F$ on $\Gamma_N$ which is exactly \eqref{eq:strong_classical_neumann} since both, $t(u)$ and $F$ belong to $L^2(\Gamma_N;\R^d)$. We continue by showing \eqref{eq:strong_contact_tangential}. Plugging \eqref{eq:decompint} in to \eqref{eq:hybineq} and using $t(u) = F$ on $\Gamma_N$ as well as $(\gamma v)|_{\Gamma_D} = 0$ for $v\in K$, we have
	\begin{equation}\label{eq:ge0}
		\int_{\Gamma_C}t(u)\cdot(\gamma v)\;\ds \ge 0\qquad\forall\,v\in K.
	\end{equation}
	Let $g_t\in H_{00}^{1/2}(\Gamma_C;\R^d)$ with $g_t\cdot\mathrm{n} = 0$, i.e.\ $g_t\in W_{C,t}$ and choose $v\in H^1(\Omega;\R^d)$ with
	\begin{equation*}
		(\gamma v)|_{\Gamma_C} = g_t,\qquad (\gamma v)|_{\Gamma_N} = 0,\qquad (\gamma v)|_{\Gamma_D} = 0.
	\end{equation*}
	Again, $\pm v\in K$ and therefore
	\begin{equation*}
		\int_{\Gamma_C}t(u)\cdot g_t\;\ds\ge 0,\qquad\text{and}\qquad -\int_{\Gamma_C}t(u)\cdot g_t\;\ds\ge 0.
	\end{equation*}
	Splitting $t(u) = \sigma_\mathrm{n}(u)\mathrm{n}+\sigma_t(u)$ into its normal and tangential components, this means
	\begin{equation*}
		\int_{\Gamma_C} (\sigma_\mathrm{n}(u)\mathrm{n}+\sigma_t(u))\cdot g_t\;\ds = \int_{\Gamma_C}\sigma_t(u)\cdot g_t\;\ds + \int_{\Gamma_C} \sigma_\mathrm{n}(u)(\mathrm{n}\cdot g_t)\;\ds = 0.\qquad\forall\,g_t\in W_{C,t}
	\end{equation*}
	With  $g_t\cdot \mathrm{n} = 0$, we thus have
	\begin{equation*}
		\int_{\Gamma_C}\sigma_t(u)\cdot g_t\;\ds = 0\qquad \forall\, g_t\in W_{C,t}
	\end{equation*}
	which means $\sigma_t(u)=0$ on $\Gamma_C$, i.e\ \eqref{eq:strong_contact_tangential}. Next, we prove \eqref{eq:strong_contact_sigma_n}. Let $\eta\in H_{00}^{1/2}(\Gamma_C)$ with $\eta\le 0$ on $\Gamma_C$, i.e. $\eta\in W_{C,\mathrm{n}}^-$ and choose $v\in H^1(\Omega;\R^d)$ with
	\begin{equation*}
		(\gamma v)|_{\Gamma_C} = \eta\,\mathrm{n},\qquad (\gamma v)|_{\Gamma_N}=0,\qquad (\gamma v)|_{\Gamma_D} = 0.
	\end{equation*}
	Then, $v_\mathrm{n} = (\gamma v)\cdot\mathrm{n} = \eta \le 0$ on $\Gamma_C$, i.e.\ $v\in K$ and thus, plugging $v$ into \eqref{eq:ge0}, we have
	\begin{equation*}
		\int_{\Gamma_C}t(u)\cdot(\eta\,\mathrm{n})\;\ds\ge 0\qquad\forall\,\eta\in W_{C,\mathrm{n}}^-.
	\end{equation*}
	Again, by splitting $t(u) = \sigma_\mathrm{n}(u)\mathrm{n}+\sigma_t(u)$ and using $\sigma_t(u) = 0$ on $\Gamma_C$, this yields
	\begin{equation*}
		\int_{\Gamma_C}\sigma_\mathrm{n}(u)\,\eta\;\ds \ge 0\qquad\forall\,\eta\in W_{C,\mathrm{n}}^-
	\end{equation*}
	or, in other words,
	\begin{equation*}
		\int_{\Gamma_C}\sigma_\mathrm{n}(u)\,\eta\;\ds \le 0\qquad\forall\,\varphi\in H_{00}^{1/2}(\Gamma_C),\;\varphi\ge 0,
	\end{equation*}
	which implies $\sigma_\mathrm{n}(u)\le 0$ on $\Gamma_C$, i.e.\ \eqref{eq:strong_contact_sigma_n}. It remains to show the complementarity condition \eqref{eq:strong_contact_comp}. Plugging \eqref{eq:decompint} for $v = u$ into \eqref{eq:hybeq} and using $t(u) = F$ on $\Gamma_N$ as well as the Dirichlet condition $(\gamma u)|_{\Gamma_D} = 0$, yields
	\begin{equation*}
		\int_{\Gamma_C}t(u)\cdot(\gamma u)\;\ds = \int_{\Gamma_C} \sigma_\mathrm{n}(u)\, u_{\mathrm{n}}\;\ds  =  0,
	\end{equation*}
	with $t(u) = \sigma_\mathrm{n}(u)\mathrm{n}$ as previously. The integrand $\sigma_\mathrm{n}(u)\, u_{\mathrm{n}}$ is nonnegative on $\Gamma_C$ since $u_\mathrm{n}\le 0$ on $\Gamma_C$ and $\sigma_\mathrm{n}(u)\le 0$ on $\Gamma_C$. Therefore $\sigma_\mathrm{n}(u)\, u_{\mathrm{n}} = 0$ on $\Gamma_C$ which is exactly \eqref{eq:strong_contact_comp}. This completes the proof.
\end{proof}
